\documentclass[11pt,letterpaper]{article}
\usepackage[utf8]{inputenc}
\usepackage[spanish]{babel}
\usepackage[margin=2.5cm]{geometry}
\usepackage{fancyhdr}
\usepackage{graphicx}
\usepackage{enumitem}
\usepackage{xcolor}
\usepackage{titlesec}
\usepackage{hyperref}

% Configuración de colores
\definecolor{wasaffblue}{RGB}{0,102,204}
\definecolor{sectiongray}{RGB}{100,100,100}

% Headers y footers
\pagestyle{fancy}
\fancyhf{}
\fancyhead[L]{\small\textcolor{wasaffblue}{WASAFF CONSULTING SpA}}
\fancyhead[R]{\small\textcolor{wasaffblue}{LEY CERO SpA}}
\fancyfoot[C]{\thepage}
\renewcommand{\headrulewidth}{0.5pt}
\renewcommand{\footrulewidth}{0.5pt}

% Formato de secciones
\titleformat{\section}
  {\normalfont\Large\bfseries\color{wasaffblue}}
  {\thesection.}{1em}{}[\titlerule]

\titleformat{\subsection}
  {\normalfont\large\bfseries\color{sectiongray}}
  {\thesubsection}{1em}{}

% Configuración de hyperlinks
\hypersetup{
    colorlinks=true,
    linkcolor=wasaffblue,
    urlcolor=wasaffblue,
    pdftitle={Contrato Wasaff-Ley Cero},
    pdfauthor={Felipe Guzmán - Nilton}
}

\begin{document}

% ====================================================================
% PORTADA
% ====================================================================
\begin{titlepage}
\centering
\vspace*{2cm}

{\Huge\bfseries\color{wasaffblue} CONTRATO DE COLABORACIÓN\\[0.5cm] TÉCNICA Y COMERCIAL\par}

\vspace{1.5cm}

{\Large\textbf{SISTEMAS DE REFRIGERACIÓN INDUSTRIAL}\par}

\vspace{3cm}

\begin{minipage}{0.45\textwidth}
\centering
\textbf{\Large WASAFF CONSULTING SpA}\\[0.3cm]
\textit{Simulación y Análisis Térmico}\\[0.2cm]
\small RUT: [COMPLETAR]\\
Rep. Legal: Felipe Andrés Guzmán Montero\\
RUT: [COMPLETAR]
\end{minipage}%
\hfill
\begin{minipage}{0.45\textwidth}
\centering
\textbf{\Large LEY CERO SpA}\\[0.3cm]
\textit{Instrumentación y Automatización}\\[0.2cm]
\small RUT: [COMPLETAR]\\
Rep. Legal: [Nilton Nombre Completo]\\
RUT: [COMPLETAR]
\end{minipage}

\vfill

{\large Santiago de Chile\\[0.2cm] \today\par}

\end{titlepage}

% ====================================================================
% CONTENIDO DEL CONTRATO
% ====================================================================

\section*{COMPARECIENTES}

En Santiago de Chile, a \rule{2cm}{0.15mm} de \rule{3cm}{0.15mm} de 2026, entre:

\vspace{0.5cm}

\textbf{DE UNA PARTE:} \textbf{FELIPE ANDRÉS GUZMÁN MONTERO}, cédula nacional de identidad número \rule{3cm}{0.15mm}, en representación de \textbf{WASAFF CONSULTING SpA}, RUT \rule{3cm}{0.15mm}, con domicilio en \rule{6cm}{0.15mm}, comuna de \rule{4cm}{0.15mm}, Santiago, en adelante denominada \textbf{"WASAFF"}.

\vspace{0.3cm}

\textbf{DE OTRA PARTE:} \textbf{[NOMBRE COMPLETO NILTON]}, cédula nacional de identidad número \rule{3cm}{0.15mm}, en representación de \textbf{LEY CERO SpA}, RUT \rule{3cm}{0.15mm}, con domicilio en \rule{6cm}{0.15mm}, comuna de \rule{4cm}{0.15mm}, Santiago, en adelante denominada \textbf{"LEY CERO"}.

\vspace{0.3cm}

Ambas partes, en adelante denominadas conjuntamente \textbf{"LAS PARTES"}, individual e indistintamente \textbf{"LA PARTE"}, han acordado celebrar el presente \textbf{CONTRATO DE COLABORACIÓN TÉCNICA Y COMERCIAL}, que se regirá por las siguientes cláusulas y condiciones:

% ====================================================================

\section{OBJETO DEL CONTRATO}

\subsection{Servicios Conjuntos}

LAS PARTES acuerdan colaborar de manera sinérgica en la prestación de los siguientes servicios en el sector de refrigeración industrial:

\begin{enumerate}[label=\alph*)]
    \item \textbf{Diagnóstico Energético:} Análisis termodinámico de sistemas de refrigeración mediante simulación computacional y validación en terreno.
    
    \item \textbf{Monitoreo Predictivo:} Implementación de sistemas de adquisición de datos (DAQ) e instrumentación IoT para monitoreo continuo de variables operacionales.
    
    \item \textbf{Gemelos Digitales:} Desarrollo de modelos digitales que replican el comportamiento térmico de instalaciones frigoríficas para optimización operacional.
    
    \item \textbf{Auditorías Térmicas:} Identificación de fugas energéticas, cálculo de eficiencias isoentrópicas y estimación de retornos de inversión (ROI) en proyectos de eficiencia energética.
\end{enumerate}

\subsection{Alcance Territorial}

Los servicios se prestarán inicialmente en la Región Metropolitana, con posibilidad de expansión a nivel nacional según demanda del mercado.

% ====================================================================

\section{ROLES Y RESPONSABILIDADES}

\subsection{Responsabilidades de WASAFF}

WASAFF se compromete a aportar las siguientes capacidades técnicas:

\begin{enumerate}[label=\arabic*.]
    \item \textbf{Modelamiento Termodinámico:} Desarrollo de simulaciones computacionales utilizando propiedades termodinámicas reales de refrigerantes mediante software especializado (CoolProp, Python).
    
    \item \textbf{Análisis de Datos:} Procesamiento de series temporales de variables operacionales (presión, temperatura, consumo eléctrico) mediante algoritmos de análisis numérico.
    
    \item \textbf{Reportería Técnica:} Generación de informes ejecutivos con diagnósticos, recomendaciones técnicas y estimaciones de impacto económico.
    
    \item \textbf{Propiedad Intelectual:} Desarrollo y mantención de algoritmos, modelos matemáticos y metodologías de análisis como activos propios de WASAFF.
    
    \item \textbf{Soporte Remoto:} Interpretación de datos, ajuste de modelos y resolución de consultas técnicas durante la ejecución de proyectos.
\end{enumerate}

\subsection{Responsabilidades de LEY CERO}

LEY CERO se compromete a aportar las siguientes capacidades operacionales:

\begin{enumerate}[label=\arabic*.]
    \item \textbf{Gestión Comercial:} Prospección de clientes, presentación de propuestas técnico-comerciales y cierre de contratos con clientes finales.
    
    \item \textbf{Instrumentación en Terreno:} Instalación, configuración y calibración de sensores de presión, temperatura y consumo eléctrico según especificaciones técnicas definidas por WASAFF.
    
    \item \textbf{Validación de Resultados:} Verificación en sitio de las predicciones del modelo, incluyendo mediciones comparativas y ajustes de instrumentación.
    
    \item \textbf{Soporte Post-Venta:} Atención de requerimientos de clientes en terreno, incluyendo mantención de equipos de medición y resolución de incidencias operacionales.
    
    \item \textbf{Logística:} Coordinación de visitas a terreno, gestión de permisos de acceso a instalaciones industriales y transporte de equipamiento.
\end{enumerate}

\subsection{Responsabilidades Compartidas}

Ambas PARTES se comprometen a:

\begin{enumerate}[label=\arabic*.]
    \item Mantener comunicación fluida y transparente durante la ejecución de proyectos.
    \item Participar conjuntamente en reuniones con clientes cuando la complejidad técnica lo requiera.
    \item Respetar los compromisos de calidad y plazos de entrega acordados con clientes.
    \item Proteger la confidencialidad de información sensible de clientes y de la contraparte.
\end{enumerate}

% ====================================================================

\section{MODELO ECONÓMICO}

\subsection{Distribución de Ingresos}

LAS PARTES acuerdan la siguiente distribución de ingresos para \textbf{TODOS} los proyectos desarrollados bajo este contrato:

\begin{center}
\fbox{\parbox{0.8\textwidth}{\centering
\textbf{\Large 50\% WASAFF \hspace{1cm} 50\% LEY CERO}\\[0.3cm]
\textit{División equitativa para todos los tipos de servicio}
}}
\end{center}

\vspace{0.3cm}

Este porcentaje se aplica sobre la \textbf{facturación neta} (monto total facturado al cliente, descontado el IVA).

\subsection{Base de Cálculo}

Los porcentajes se calcularán sobre el \textbf{ingreso neto} de cada proyecto, definido como:

\begin{equation}
\text{Ingreso Neto} = \text{Facturación Total} - \text{IVA} - \text{Costos Directos Compartidos}
\end{equation}

Donde \textbf{Costos Directos Compartidos} incluyen:
\begin{itemize}
    \item Viáticos y gastos de traslado a terreno
    \item Arriendo de equipamiento especializado no disponible por LAS PARTES
    \item Materiales consumibles requeridos para el proyecto específico
    \item Subcontratación de servicios especializados (termografía, análisis de aceite, etc.)
\end{itemize}

Los costos operacionales propios de cada PARTE (oficinas, equipos propios, seguros, etc.) \textbf{NO} se descuentan de la base de cálculo.

\subsection{Facturación}

\begin{enumerate}[label=\arabic*.]
    \item \textbf{Modalidad Preferida:} Cada PARTE factura directamente al cliente su porcentaje correspondiente (50\% cada una).
    
    \item \textbf{Modalidad Alternativa:} Si el cliente requiere factura única, una PARTE factura el 100\% y liquida a la otra dentro de \textbf{5 días corridos} desde la recepción del pago del cliente.
    
    \item La PARTE que facture el 100\% deberá enviar a la otra:
    \begin{itemize}
        \item Copia de factura emitida al cliente
        \item Comprobante de pago recibido del cliente
        \item Detalle de costos directos compartidos (si aplica)
        \item Transferencia bancaria del 50\% correspondiente
    \end{itemize}
\end{enumerate}

\subsection{Transparencia Financiera}

Cada proyecto tendrá una hoja de cálculo compartida (Google Sheets o Excel) donde se registrarán:
\begin{itemize}
    \item Monto cotizado al cliente
    \item Costos directos compartidos (con respaldos)
    \item Base de cálculo resultante
    \item Monto correspondiente a cada PARTE
\end{itemize}

% ====================================================================

\section{PROPIEDAD INTELECTUAL}

\subsection{Activos de WASAFF}

Son propiedad exclusiva de WASAFF y no se transfieren ni licencian a LEY CERO:

\begin{enumerate}[label=\alph*)]
    \item Código fuente de algoritmos y scripts de simulación
    \item Modelos matemáticos y físicos desarrollados
    \item Metodologías de análisis termodinámico
    \item Bibliotecas de propiedades de refrigerantes
    \item Plantillas de reportes técnicos (estructura y contenido)
\end{enumerate}

\subsection{Activos de LEY CERO}

Son propiedad exclusiva de LEY CERO y no se transfieren a WASAFF:

\begin{enumerate}[label=\alph*)]
    \item Base de datos de clientes y contactos comerciales
    \item Protocolos de instalación y calibración de instrumentación
    \item Diseños de hardware y sistemas SCADA propietarios
    \item Procedimientos operacionales estándar (SOP) de LEY CERO
\end{enumerate}

\subsection{Propiedad Compartida}

Los siguientes elementos son propiedad compartida de ambas PARTES:

\begin{enumerate}[label=\alph*)]
    \item Reportes técnicos entregados a clientes (cada PARTE puede incluirlos en su portafolio con autorización de la otra)
    \item Presentaciones comerciales conjuntas
    \item Material de difusión (casos de éxito, testimoniales)
\end{enumerate}

\subsection{Uso Post-Término}

En caso de término del contrato, cada PARTE conserva:
\begin{itemize}
    \item Derecho a mostrar proyectos realizados en conjunto como referencia comercial
    \item Prohibición de replicar metodologías o código fuente de la otra PARTE
    \item Derecho a contactar clientes comunes respetando cláusula de no competencia
\end{itemize}

% ====================================================================

\section{CONFIDENCIALIDAD}

\subsection{Información Confidencial}

Cada PARTE se compromete a mantener en estricta confidencialidad:

\begin{enumerate}[label=\arabic*.]
    \item Datos operacionales de clientes (temperaturas, presiones, consumos, productos almacenados)
    \item Información financiera de proyectos y propuestas comerciales
    \item Metodologías técnicas y algoritmos de la contraparte
    \item Estrategias comerciales y pricing
\end{enumerate}

\subsection{Excepciones}

No se considera confidencial la información que:
\begin{itemize}
    \item Sea de dominio público
    \item Deba revelarse por mandato legal o judicial
    \item Sea autorizada por escrito por la PARTE afectada
\end{itemize}

\subsection{Vigencia}

Esta obligación de confidencialidad se mantiene \textbf{indefinidamente}, incluso después del término del presente contrato.

% ====================================================================

\section{EXCLUSIVIDAD Y NO COMPETENCIA}

\subsection{Sin Exclusividad General}

Este contrato \textbf{NO es exclusivo}. Ambas PARTES pueden:
\begin{itemize}
    \item Trabajar con otros partners en proyectos no relacionados con refrigeración industrial
    \item Ofrecer servicios independientes en otros sectores industriales
    \item Participar en licitaciones públicas de manera individual
\end{itemize}

\subsection{Protección de Clientes Conjuntos}

Durante la vigencia del contrato y \textbf{12 meses posteriores a su término}, ninguna PARTE podrá:

\begin{enumerate}[label=\arabic*.]
    \item Ofrecer servicios de diagnóstico energético + instrumentación a clientes prospectados por la otra PARTE sin autorización escrita.
    
    \item Utilizar información técnica obtenida durante la colaboración para competir directamente con proyectos similares.
    
    \item Contratar personal clave de la otra PARTE con el objetivo de replicar las capacidades técnicas del presente contrato.
\end{enumerate}

\textbf{Clientes conjuntos:} Se definen como aquellos donde ambas PARTES participaron en la prospección, cotización o ejecución del proyecto.

% ====================================================================

\section{VIGENCIA Y TÉRMINO}

\subsection{Vigencia Inicial}

El presente contrato tendrá una vigencia de \textbf{24 meses} contados desde la fecha de firma, con renovación automática por períodos de 12 meses salvo aviso contrario.

\subsection{Término Anticipado}

Cualquiera de LAS PARTES puede poner término al contrato en las siguientes circunstancias:

\begin{enumerate}[label=\alph*)]
    \item \textbf{Término Unilateral:} Con 60 días de aviso previo por escrito, sin expresión de causa.
    
    \item \textbf{Término por Incumplimiento Grave:} Inmediato, ante situaciones como:
    \begin{itemize}
        \item Incumplimiento reiterado de obligaciones contractuales
        \item Violación de confidencialidad
        \item Fraude o representación falsa de capacidades técnicas
        \item Conflictos de interés no declarados
    \end{itemize}
    
    \item \textbf{Término por Fuerza Mayor:} Ante eventos que imposibiliten la continuidad (desastres naturales, cambios regulatorios, quiebra de alguna PARTE).
\end{enumerate}

\subsection{Efectos del Término}

Al término del contrato:

\begin{enumerate}[label=\arabic*.]
    \item Proyectos en curso se completarán bajo las condiciones originales del contrato.
    \item Se realizará liquidación final de pagos pendientes dentro de 30 días.
    \item Cada PARTE conserva su propiedad intelectual según lo establecido en la Cláusula Cuarta.
    \item Se mantiene la obligación de confidencialidad indefinidamente.
    \item Inicia el período de no competencia de 12 meses para clientes conjuntos.
\end{enumerate}

% ====================================================================

\section{RESOLUCIÓN DE CONFLICTOS}

\subsection{Procedimiento Escalonado}

Ante cualquier controversia relacionada con la interpretación o cumplimiento del presente contrato, LAS PARTES se comprometen a:

\begin{enumerate}[label=\textbf{Etapa \arabic*:}]
    \item \textbf{Negociación Directa} (15 días hábiles): Los representantes legales de ambas PARTES se reunirán para buscar una solución de común acuerdo.
    
    \item \textbf{Mediación} (30 días hábiles): Si no hay acuerdo, se someterán a mediación ante un mediador inscrito en el Registro de Mediadores del Ministerio de Justicia, cuyos honorarios se dividirán 50\%-50\%.
    
    \item \textbf{Arbitraje} (último recurso): Si la mediación fracasa, se someterán a arbitraje de la Cámara de Comercio de Santiago (CAM), con árbitro mixto en cuanto al procedimiento y arbitrador en cuanto al fallo.
\end{enumerate}

\subsection{Jurisdicción}

Para todos los efectos legales derivados del presente contrato, LAS PARTES fijan domicilio en la ciudad y comuna de Santiago, sometiéndose a la competencia de sus Tribunales Ordinarios de Justicia, renunciando expresamente a cualquier otro fuero.

% ====================================================================

\section{DISPOSICIONES GENERALES}

\subsection{Integridad del Contrato}

Este documento constituye el acuerdo completo entre LAS PARTES, dejando sin efecto cualquier negociación, acuerdo verbal o correspondencia previa.

\subsection{Modificaciones}

Cualquier modificación al presente contrato deberá constar por escrito y ser firmada por los representantes legales de ambas PARTES.

\subsection{Cesión}

Ninguna PARTE podrá ceder o transferir sus derechos y obligaciones bajo este contrato sin el consentimiento previo y por escrito de la otra PARTE.

\subsection{Notificaciones}

Todas las comunicaciones formales se realizarán por correo electrónico a:

\begin{itemize}
    \item WASAFF: contacto@wasaff.cl
    \item LEY CERO: [correo de Nilton]
\end{itemize}

Con copia física certificada al domicilio indicado en el encabezado del presente contrato.

\subsection{Independencia de las PARTES}

Este contrato no constituye una sociedad, asociación o joint venture. Cada PARTE mantiene su independencia legal, financiera y operacional.

% ====================================================================

\vspace{1cm}

En comprobante y aceptación de todo lo anterior, las PARTES firman el presente contrato en duplicado, quedando un ejemplar en poder de cada una de ellas, en Santiago de Chile, a \rule{2cm}{0.15mm} de \rule{3cm}{0.15mm} de 2026.

\vspace{2cm}

\begin{minipage}{0.45\textwidth}
\centering
\rule{6cm}{0.4pt}\\
\textbf{FELIPE ANDRÉS GUZMÁN MONTERO}\\
WASAFF CONSULTING SpA\\
RUT: [COMPLETAR]\\
Rep. Legal
\end{minipage}%
\hfill
\begin{minipage}{0.45\textwidth}
\centering
\rule{6cm}{0.4pt}\\
\textbf{[NOMBRE COMPLETO NILTON]}\\
LEY CERO SpA\\
RUT: [COMPLETAR]\\
Rep. Legal
\end{minipage}

\end{document}
